\section{Introduction}
Superconductivity with nonzero center-of-mass pairing momentum provides a route beyond a spatially uniform condensate. Fulde and Ferrell and, independently, Larkin and Ovchinnikov established the finite-momentum pairing and spatially modulated mechanisms now grouped under the FFLO setting~\cite{FuldeFerrell1964,LarkinOvchinnikov1965}. Subsequent reviews have developed the thermodynamic, material, and ultracold-gas consequences of Fermi-surface mismatch and competing pairing momenta~\cite{CasalbuoniNardulli2004,MatsudaShimahara2007,KinnunenEtAl2018}. The central phase-selection problem is therefore intrinsically thermodynamic: competing minima must be compared over both pairing amplitude and momentum.

The search for FFLO structure spans several physical platforms. Heavy-fermion and organic superconductors have displayed field-induced phases and spectroscopic, calorimetric, or magnetic signatures discussed in relation to FFLO order~\cite{RadovanEtAl2003,BianchiEtAl2003,KenzelmannEtAl2008,LortzEtAl2007,MayaffreEtAl2014}; these representative studies motivate the problem without implying that every identification is universal or uncontested. Population-imbalanced ultracold Fermi gases provide a complementary tunable setting for pairing under mismatch, phase separation, and finite-momentum phases~\cite{LiaoEtAl2010,ZwierleinEtAl2006,PartridgeEtAl2006,SheehyRadzihovsky2007,RadzihovskySheehy2010}.

Altermagnets add a distinct route to this competition. Their compensated magnetic order can nevertheless generate symmetry-structured momentum-dependent spin splitting~\cite{SmejkalSinovaJungwirth2022,SmejkalEtAl2022,Mazin2022,SmejkalEtAl2020}. Model studies have connected this splitting to zero-field finite-momentum or other unconventional pairing states~\cite{ChakrabortyBlackSchaffer2024,HongParkKim2025,BoseVadnaisParamekanti2024}, and neighboring altermagnetic constructions have been proposed for topological superconductivity and Majorana modes~\cite{ZhuEtAl2023,GhorashiHughesCano2024}. These precedents motivate an altermagnetic finite-$q$ problem but do not determine the phase structure of the one-dimensional model studied here.

Topology adds a second, logically separate requirement. Finite-momentum pairing can coexist with topological superfluidity in spin-orbit-coupled models~\cite{ZhangYi2013,QuEtAl2013,CaoEtAl2014}. More generally, bulk symmetries and invariants organize topological superconducting phases~\cite{ReadGreen2000,SatoAndo2017,QiZhang2011,SchnyderEtAl2008,RyuEtAl2010}. Majorana-platform studies connect those bulk distinctions to boundary physics~\cite{LutchynSauDasSarma2010,OregRefaelVonOppen2010,Alicea2012,Beenakker2013}. In a gapped one-dimensional Bogoliubov--de Gennes (BdG) system in symmetry class D, the relevant $\mathbb Z_2$ distinction can be expressed through a Pfaffian invariant~\cite{Kitaev2001,BudichArdonne2013}; its use is meaningful only after bulk-gap applicability is established. A topology-resolved FFLO map must therefore identify a reliable thermodynamic minimum before classifying trusted gapped FFLO states.

This hierarchy is computationally demanding because every accepted parameter point requires minimization of $F(\Delta,q)$, comparison with the normal reference, adaptive checks of shallow or edge minima, and, where applicable, full-zone spectral and Pfaffian diagnostics. Active learning and Bayesian experimental design formalize the selection of informative sequential evaluations~\cite{CohnAtlasLadner1994,MacKay1992,ShahriariEtAl2016}; phase-diagram applications demonstrate boundary-focused sampling~\cite{TerayamaEtAl2019,DaiBrussGlotzer2020,LookmanEtAl2019}, while adaptive sampling has also reduced expensive electronic-structure, atomistic, and chemical search budgets~\cite{TranUlissi2018,PodryabinkinShapeev2017,JinnouchiEtAl2019,VandermauseEtAl2020,BatraEtAl2021,NoeEtAl2019}. For a physical phase diagram, however, an acquisition model should decide where to calculate next, not replace the exact calculation that defines the phase.

Here we develop a reliability-aware active-learning workflow in which an audited exact BdG oracle remains the sole authority for thermodynamic and topological labels. We first reconstruct the two-dimensional Normal, uniform-superconducting, and FFLO structure while retaining a separate hard-risk numerical layer. On trusted gapped FFLO states, full-zone gap and Pfaffian diagnostics resolve cFFLO/tFFLO structure with an explicit neighborhood-size sensitivity audit. Finally, we extend the same oracle--scheduler separation to a three-dimensional parameter domain. The extension reveals thermodynamic and topology-aware structure, while its strict topology convergence criteria remain open; it is not presented as a complete three-dimensional topology atlas.
